\section{Circuit Complexity --- Hard Wire}
%%% generated by the blueprint pipeline from TCSlib.Complexity.CircuitComplexity.HardWire
\subsection*{Circuit Complexity --- Hard Wire: Introduction}

Declarations of the Lean module \texttt{TCSlib.Complexity.CircuitComplexity.HardWire}.

\subsection*{Circuit Complexity --- Hard Wire: Results}

\begin{theorem}[Evaluation of constant gates]
\lean{BoolCircuit.runWith_eval_constGates}
\label{BoolCircuit.runWith_eval_constGates}
\leanok
\uses{BoolCircuit.DAGGate.eval, BoolCircuit.constGate, BoolCircuit.runWith, BoolCircuit.runWith_cons}
\difficulty{4}
For any Boolean lists $l$ and $v$, successively evaluating the constant gates
corresponding to the entries of $l$ and appending their outputs to the initial list $v$
yields exactly $v$ followed by $l$.
\end{theorem}

\begin{theorem}[Depths of constant gates]
\lean{BoolCircuit.runWith_depth_constGates}
\label{BoolCircuit.runWith_depth_constGates}
\leanok
\uses{BoolCircuit.DAGGate.depth, BoolCircuit.constGate, BoolCircuit.runWith, BoolCircuit.runWith_cons}
\difficulty{4}
Let $l$ be a list of Boolean values and let $I$ be a list of natural numbers. Starting
with the values in $I$, append the depth of each constant gate associated with an entry
of $l$, in the order of the entries. The resulting list is $I$ followed by $l.length$
copies of $1$.
\end{theorem}

\begin{theorem}[Bound on maxima under pointwise slack]
\lean{BoolCircuit.foldr_max_le_succ}
\label{BoolCircuit.foldr_max_le_succ}
\leanok
\difficulty{5}
Let $l$ be a finite list of natural numbers and let $p,q:\bbr	o\bbr$ be functions on the
natural numbers. If $p(a)\le q(a)+1$ for every $a$ in $l$, then the maximum of the
values $p(a)$ over $a$ in $l$, taken to be $0$ when $l$ is empty, is at most one plus
the corresponding maximum of the values $q(a)$.
\end{theorem}

\begin{theorem}[Depth under remapping]
\lean{BoolCircuit.DAGGate.depth_remap_le}
\label{BoolCircuit.DAGGate.depth_remap_le}
\leanok
\uses{BoolCircuit.DAGGate, BoolCircuit.DAGGate.depth, BoolCircuit.DAGGate.remap, BoolCircuit.foldr_max_le_succ}
\difficulty{3}
Let $g$ be a Boolean circuit gate, let $ds,ds'$ be lists of natural numbers, and let
$\sigma$ be a map from natural numbers to natural numbers. Suppose that for every argument
$a$ of $g$, the defaulted lookup in $ds$ at index $\sigma(a)$ is at most the defaulted
lookup in $ds'$ at index $a$, plus $1$. Then the depth of the remapped gate computed
from $ds$ is at most the depth of $g$ computed from $ds'$, plus $1$.
\end{theorem}

\begin{theorem}[Bounded Default List Lookup]
\lean{BoolCircuit.getD_le_one}
\label{BoolCircuit.getD_le_one}
\leanok
\difficulty{4}
Let $l$ be a list of natural numbers whose entries are all at most $1$. For every
natural number $k$, the value at index $k$, using $0$ as the default when the index is
out of range, is at most $1$.
\end{theorem}

\begin{definition}[Hardwiring circuit inputs]
\lean{BoolCircuit.DAGCircuit.hardwire}
\label{BoolCircuit.DAGCircuit.hardwire}
\leanok
\uses{BoolCircuit.DAGCircuit, BoolCircuit.constGate}
Let $C$ be a directed acyclic Boolean circuit with $n+m$ input vertices, and let $y$ be
a Boolean assignment to the last $m$ inputs. Form a circuit with $n$ input vertices by
placing, before the gates of $C$, one constant gate for each value of $y$, in order;
retain the gates of $C$ and its output vertex.
\end{definition}

\begin{theorem}[Values after hardwiring]
\lean{BoolCircuit.DAGCircuit.hardwire_values}
\label{BoolCircuit.DAGCircuit.hardwire_values}
\leanok
\uses{BoolCircuit.DAGCircuit.hardwire, BoolCircuit.DAGCircuit.values, BoolCircuit.runWith_append, BoolCircuit.runWith_eval_constGates}
\difficulty{1}
Let $C$ be a Boolean circuit on $n+m$ inputs, let $y$ assign Boolean values to the last
$m$ inputs, and let $x$ assign Boolean values to the first $n$ inputs. The list of
vertex values produced by the circuit obtained by hardwiring $y$ at $x$ equals the list
produced by $C$ at the concatenated assignment $(x,y)$.
\end{theorem}

\begin{theorem}[Evaluation after hardwiring inputs]
\lean{BoolCircuit.DAGCircuit.hardwire_eval}
\label{BoolCircuit.DAGCircuit.hardwire_eval}
\leanok
\uses{BoolCircuit.DAGCircuit.eval, BoolCircuit.DAGCircuit.hardwire, BoolCircuit.DAGCircuit.hardwire_values}
\difficulty{3}
Let $C$ be a Boolean circuit with $n+m$ input vertices, let $y$ assign Boolean values to
the last $m$ inputs, and let $x$ assign values to the first $n$. Evaluating the circuit
obtained by hardwiring $y$ gives the same output as evaluating $C$ on the concatenated
assignment $(x,y)$.
\end{theorem}

\begin{theorem}[Size after hardwiring]
\lean{BoolCircuit.DAGCircuit.hardwire_size}
\label{BoolCircuit.DAGCircuit.hardwire_size}
\leanok
\uses{BoolCircuit.DAGCircuit.hardwire, BoolCircuit.DAGCircuit.size}
\difficulty{3}
Let $C$ be a Boolean circuit on $n+m$ inputs and let $y$ assign Boolean values to the
last $m$ inputs. The size of the circuit obtained by hardwiring $y$ equals the size of
$C$.
\end{theorem}

\begin{theorem}[Size bound after hardwiring]
\lean{BoolCircuit.DAGCircuit.hardwire_size_le}
\label{BoolCircuit.DAGCircuit.hardwire_size_le}
\leanok
\uses{BoolCircuit.DAGCircuit.hardwire, BoolCircuit.DAGCircuit.hardwire_size, BoolCircuit.DAGCircuit.size}
\difficulty{1}
Let $C$ be a Boolean circuit on $n+m$ inputs and let $y$ assign Boolean values to the
last $m$ inputs. The size of the circuit obtained by hardwiring $y$ is at most the size
of $C$.
\end{theorem}

\begin{theorem}[Preservation of well-formedness]
\lean{BoolCircuit.DAGCircuit.hardwire_isWellFormed}
\label{BoolCircuit.DAGCircuit.hardwire_isWellFormed}
\leanok
\uses{BoolCircuit.DAGCircuit.IsWellFormed, BoolCircuit.DAGCircuit.hardwire, BoolCircuit.constGate_kind_ne_not}
\difficulty{4}
Let $C$ be a Boolean circuit on $n+m$ inputs and let $y$ assign Boolean values to the
last $m$ inputs. If every gate of $C$ has distinct input vertices and each NOT gate has
exactly one input, then every gate of the circuit obtained by hardwiring $y$ also has
distinct input vertices, and each of its NOT gates has exactly one input.
\end{theorem}

\begin{theorem}[Preservation of fan-in at most two]
\lean{BoolCircuit.DAGCircuit.hardwire_isFaninTwo}
\label{BoolCircuit.DAGCircuit.hardwire_isFaninTwo}
\leanok
\uses{BoolCircuit.DAGCircuit.IsFaninTwo, BoolCircuit.DAGCircuit.hardwire, BoolCircuit.DAGCircuit.hardwire_isWellFormed}
\difficulty{4}
Let $C$ be a Boolean circuit on $n+m$ inputs and let $y$ assign Boolean values to the
last $m$ inputs. If every gate of $C$ has distinct input vertices, every NOT gate has
exactly one input, and every gate has at most two inputs, then the circuit obtained by
hardwiring $y$ has these same properties.
\end{theorem}

\begin{theorem}[Depth after hardwiring inputs]
\lean{BoolCircuit.DAGCircuit.hardwire_depth_le}
\label{BoolCircuit.DAGCircuit.hardwire_depth_le}
\leanok
\uses{BoolCircuit.DAGCircuit.depth, BoolCircuit.DAGCircuit.depthAt, BoolCircuit.DAGCircuit.depths, BoolCircuit.DAGCircuit.hardwire, BoolCircuit.DAGGate.depth, BoolCircuit.DAGGate.depth_remap_le, BoolCircuit.DAGGate.remap, BoolCircuit.DAGGate.remap_id, BoolCircuit.getD_le_one, BoolCircuit.runWith_append, BoolCircuit.runWith_depth_constGates, BoolCircuit.runWith_remap_rel}
\difficulty{6}
Let $C$ be a directed acyclic circuit with $n+m$ input bits, and fix an assignment to
its last $m$ inputs. The circuit obtained by hardwiring those inputs has depth at most
$	ext{depth}(C)+1$.
\end{theorem}

\begin{definition}[Permutation of input indices]
\lean{BoolCircuit.DAGCircuit.leftPerm}
\label{BoolCircuit.DAGCircuit.leftPerm}
\leanok
For natural numbers $m,n$ and an index $v$, define a new index by sending $v$ to $n+v$
if $v<m$, to $v-m$ if $m\le v<m+n$, and leaving it equal to $v$ if $v\ge m+n$.
\end{definition}

\begin{theorem}[Bound on permuted indices]
\lean{BoolCircuit.DAGCircuit.leftPerm_lt}
\label{BoolCircuit.DAGCircuit.leftPerm_lt}
\leanok
\uses{BoolCircuit.DAGCircuit.leftPerm}
\difficulty{4}
For natural numbers $m,n,k,v$, if $v<m+n+k$, then the interval permutation that sends
$v<m$ to $n+v$, sends $m\le v<m+n$ to $v-m$, and fixes $v\ge m+n$ sends $v$ to a number
less than $n+m+k$.
\end{theorem}

\begin{theorem}[Permutation of gate indices]
\lean{BoolCircuit.DAGCircuit.leftPerm_gate}
\label{BoolCircuit.DAGCircuit.leftPerm_gate}
\leanok
\uses{BoolCircuit.DAGCircuit.leftPerm}
\difficulty{4}
For all natural numbers $m,n,i$, the index permutation that moves the first block of $m$
indices to the end sends the index $m+n+i$ to $n+m+i$.
\end{theorem}

\begin{theorem}[Injectivity of the interval permutation]
\lean{BoolCircuit.DAGCircuit.leftPerm_injective}
\label{BoolCircuit.DAGCircuit.leftPerm_injective}
\leanok
\uses{BoolCircuit.DAGCircuit.leftPerm}
\difficulty{4}
For natural numbers $m,n$, define a map on the natural numbers by sending $v<m$ to
$n+v$, sending $m\le v<m+n$ to $v-m$, and fixing $v\ge m+n$. This map is injective.
\end{theorem}

\begin{definition}[Hardwiring left inputs]
\lean{BoolCircuit.DAGCircuit.hardwireLeft}
\label{BoolCircuit.DAGCircuit.hardwireLeft}
\leanok
\uses{BoolCircuit.DAGCircuit, BoolCircuit.DAGCircuit.leftPerm, BoolCircuit.DAGCircuit.leftPerm_lt, BoolCircuit.DAGGate.remap, BoolCircuit.constGate}
Let $C$ be a directed acyclic Boolean circuit with $m+n$ input vertices, and let $y$ be
a Boolean assignment to the first $m$ inputs. Form a circuit with $n$ input vertices by
placing, before the gates of $C$, one constant gate for each value of $y$, in order.
Remap every argument index of the original gates and the output index by sending $v$ to
$n+v$ for $v<m$, to $v-m$ for $m\le v<m+n$, and to $v$ for $v\ge m+n$.
\end{definition}

\begin{theorem}[Evaluation after hardwiring left inputs]
\lean{BoolCircuit.DAGCircuit.hardwireLeft_eval}
\label{BoolCircuit.DAGCircuit.hardwireLeft_eval}
\leanok
\uses{BoolCircuit.DAGCircuit.eval, BoolCircuit.DAGCircuit.hardwireLeft, BoolCircuit.DAGCircuit.leftPerm, BoolCircuit.DAGCircuit.leftPerm_gate, BoolCircuit.DAGCircuit.leftPerm_lt, BoolCircuit.DAGCircuit.values, BoolCircuit.DAGGate.eval, BoolCircuit.DAGGate.eval_remap, BoolCircuit.runWith_append, BoolCircuit.runWith_eval_constGates, BoolCircuit.runWith_remap_rel}
\difficulty{5}
Let $C$ be a directed acyclic circuit with $m+n$ input bits, fix an assignment $y$ to
its first $m$ inputs, and let $x$ assign Boolean values to the remaining $n$ inputs. The
evaluation of the circuit obtained by hardwiring the first inputs to $y$ at $x$ equals
the evaluation of $C$ at the concatenated assignment $(y,x)$.
\end{theorem}

\begin{theorem}[Size after hardwiring left inputs]
\lean{BoolCircuit.DAGCircuit.hardwireLeft_size}
\label{BoolCircuit.DAGCircuit.hardwireLeft_size}
\leanok
\uses{BoolCircuit.DAGCircuit.hardwireLeft, BoolCircuit.DAGCircuit.size}
\difficulty{3}
Let $C$ be a directed acyclic circuit with $m+n$ input bits, and fix an assignment to
its first $m$ inputs. The size of the resulting circuit, defined as its number of input
bits plus its number of gates, equals the size of $C$.
\end{theorem}

\begin{theorem}[Size bound after hardwiring left inputs]
\lean{BoolCircuit.DAGCircuit.hardwireLeft_size_le}
\label{BoolCircuit.DAGCircuit.hardwireLeft_size_le}
\leanok
\uses{BoolCircuit.DAGCircuit.hardwireLeft, BoolCircuit.DAGCircuit.hardwireLeft_size, BoolCircuit.DAGCircuit.size}
\difficulty{1}
Let $C$ be a directed acyclic circuit with $m+n$ input bits, and fix an assignment to
its first $m$ inputs. The size of the resulting circuit, defined as its number of input
bits plus its number of gates, is at most the size of $C$.
\end{theorem}

\begin{theorem}[Preservation of well-formedness]
\lean{BoolCircuit.DAGCircuit.hardwireLeft_isWellFormed}
\label{BoolCircuit.DAGCircuit.hardwireLeft_isWellFormed}
\leanok
\uses{BoolCircuit.DAGCircuit.IsWellFormed, BoolCircuit.DAGCircuit.hardwireLeft, BoolCircuit.DAGCircuit.leftPerm_injective, BoolCircuit.DAGGate.remap, BoolCircuit.constGate_kind_ne_not}
\difficulty{4}
Let $C$ be a directed acyclic circuit with $m+n$ input bits in which each gate has
distinct arguments and every NOT gate has exactly one argument. For any assignment to
the first $m$ inputs, the circuit obtained by hardwiring those inputs also has distinct
arguments at each gate, and every NOT gate in it has exactly one argument.
\end{theorem}

\begin{theorem}[Preservation of fan-in two]
\lean{BoolCircuit.DAGCircuit.hardwireLeft_isFaninTwo}
\label{BoolCircuit.DAGCircuit.hardwireLeft_isFaninTwo}
\leanok
\uses{BoolCircuit.DAGCircuit.IsFaninTwo, BoolCircuit.DAGCircuit.hardwireLeft, BoolCircuit.DAGCircuit.hardwireLeft_isWellFormed, BoolCircuit.DAGGate.remap}
\difficulty{4}
Let $C$ be a directed acyclic circuit with $m+n$ input bits and at most two arguments at
each gate; assume also that each gate has distinct arguments and every NOT gate has
exactly one argument. For any assignment to the first $m$ inputs, the circuit obtained
by hardwiring those inputs still has at most two arguments at each gate, distinct
arguments at each gate, and exactly one argument at every NOT gate.
\end{theorem}

\begin{theorem}[Depth after hardwiring left inputs]
\lean{BoolCircuit.DAGCircuit.hardwireLeft_depth_le}
\label{BoolCircuit.DAGCircuit.hardwireLeft_depth_le}
\leanok
\uses{BoolCircuit.DAGCircuit.depth, BoolCircuit.DAGCircuit.depthAt, BoolCircuit.DAGCircuit.depths, BoolCircuit.DAGCircuit.hardwireLeft, BoolCircuit.DAGCircuit.leftPerm, BoolCircuit.DAGCircuit.leftPerm_gate, BoolCircuit.DAGCircuit.leftPerm_lt, BoolCircuit.DAGGate.depth, BoolCircuit.DAGGate.depth_remap_le, BoolCircuit.getD_le_one, BoolCircuit.runWith_append, BoolCircuit.runWith_depth_constGates, BoolCircuit.runWith_remap_rel}
\difficulty{6}
Let $C$ be a directed acyclic circuit with $m+n$ input bits, and fix an assignment $y$
to its first $m$ inputs. The circuit obtained by hardwiring those inputs has depth at
most $	ext{depth}(C)+1$.
\end{theorem}

\begin{definition}[Hardwiring a circuit family]
\lean{BoolCircuit.DAGCircuitFamily.hardwire}
\label{BoolCircuit.DAGCircuitFamily.hardwire}
\leanok
\uses{BoolCircuit.DAGCircuit, BoolCircuit.DAGCircuit.hardwire, BoolCircuit.DAGCircuitFamily}
Let $a$ assign a natural number to each input length. For each $n$, let $D_n$ be a
directed acyclic Boolean circuit with $n+a(n)$ inputs, and let $\alpha_n$ be a Boolean
assignment to its last $a(n)$ inputs. The resulting circuit family assigns to each $n$
the circuit obtained by hardwiring $\alpha_n$ into $D_n$.
\end{definition}

\begin{theorem}[Hardwired circuit language]
\lean{BoolCircuit.DAGCircuitFamily.hardwire_language}
\label{BoolCircuit.DAGCircuitFamily.hardwire_language}
\leanok
\uses{BoolCircuit.DAGCircuit.eval, BoolCircuit.DAGCircuit.hardwire_eval, BoolCircuit.DAGCircuitFamily.hardwire, BoolCircuit.DAGCircuitFamily.language, BoolCircuit.DAGCircuitFamily.mem_language_iff}
\difficulty{3}
Let $a$ assign a natural number of fixed input bits to each input length, let $D_n$ be a
Boolean circuit on $n+a(n)$ inputs for each $n$, and choose an assignment $\alpha_n$ to
those $a(n)$ bits. The language accepted by the family obtained by fixing these bits to
$\alpha_n$ consists exactly of the words $w$ for which $D_{|w|}$ accepts the
concatenation of $w$ and $\alpha_{|w|}$.
\end{theorem}

\begin{theorem}[Circuit size after hardwiring inputs]
\lean{Language.InSIZE.of_hardwire}
\label{Language.InSIZE.of_hardwire}
\leanok
\uses{BoolCircuit.DAGCircuit, BoolCircuit.DAGCircuit.IsFaninTwo, BoolCircuit.DAGCircuit.eval, BoolCircuit.DAGCircuit.hardwire_isFaninTwo, BoolCircuit.DAGCircuit.hardwire_size_le, BoolCircuit.DAGCircuit.size, BoolCircuit.DAGCircuitFamily.hardwire, BoolCircuit.DAGCircuitFamily.hardwire_language, Language.InSIZE}
\difficulty{1}
Let $a,p : \mathbb{N} \to \mathbb{N}$, and for each $n$ let $D_n$ be a Boolean circuit
on $n+a(n)$ input bits and let $\alpha_n$ be a fixed Boolean string of length $a(n)$.
Suppose every $D_n$ is well formed with fan-in at most two, and its size is at most
$p(n)$. Then the language consisting of Boolean strings $w$ for which $D_{|w|}$ accepts
the concatenation of $w$ and $\alpha_{|w|}$ is recognized by a family of fan-in-two
Boolean circuits of size at most $p(n)$ on inputs of length $n$.
\end{theorem}

\begin{theorem}[Hardwiring preserves polynomial size]
\lean{Language.InPPoly.of_hardwire}
\label{Language.InPPoly.of_hardwire}
\leanok
\uses{BoolCircuit.DAGCircuit, BoolCircuit.DAGCircuit.IsFaninTwo, BoolCircuit.DAGCircuit.eval, BoolCircuit.DAGCircuit.size, Language.InPPoly, Language.InSIZE.of_hardwire}
\difficulty{1}
Let $a : \mathbb{N} \to \mathbb{N}$, and for each input length $n$ let $D_n$ be a
Boolean circuit on $n+a(n)$ input bits. For each $n$, fix an assignment $\alpha_n$ to
the final $a(n)$ input bits. Suppose every $D_n$ is well-formed with fan-in at most two,
and there are natural numbers $c,k$ such that the size of $D_n$ is at most $c(n+1)^k$
for every $n$. Then the language of strings $w$ accepted by $D_{|w|}$ when given $w$
followed by the fixed assignment $\alpha_{|w|}$ has polynomial circuit size with fan-in
at most two.
\end{theorem}
